\section{Tradition on Italian Soil}

\begin{quotex}
From Pythagoras to Virgil to Dante, the chain of Tradition was never lost on Italian Soil.
\flright{\textsc{Rene Guenon}}

In old Europe human life received its ideal content from the Catholic faith, on the one hand, and from knightly feudalism on the other
\flright{\textsc{Vladimir Solovyov.}}
\end{quotex}


\paragraph{Pythagoras}


Pythagoras established his first community in Crotone, bringing esoteric teachings onto the Italian peninsula. From Pythagoras, there is a chain leading back to Orpheus and ultimately to Hermes Trismegistus. The relationship of this Tradition to others is described in the \textit{Our Father Course}\footnote{You can find such course in its own book of this Series.}.

\paragraph{Virgil}

Virgil continues this line, as revealed in his poem \emph{Aeneid} and prophesied in the \emph{Eclogues\footnote{\url{https://gornahoor.net/?p=5481}}}. The Western initiatic tradition was the story of the encounter with Divine Wisdom of Sophia. Aeneas, the progenitor of the founders of Rome, was a type of this tradition. Numa Pompilius, the second king of Rome often consulted the nymph Egeria who taught him wisdom.

Anticipating Dante's Journey, Aeneas descends into hell, led by the Cumaean Sibyl\footnote{\url{https://en.wikipedia.org/wiki/Cumaean_Sibyl}}. Following that, he is led to the Elysian Fields. There the righteous and the heroic lived lives of natural happiness, which was an extension to the pleasures they had enjoyed in life. This is equivalent to the Limbo of the Fathers. This is reserved to the righteous and heroic reflective of the two radically different conceptions of happiness\footnote{\url{https://gornahoor.net/?p=13}}.

As pleasant as that may be, it ultimately falls short of deliverance, or better said, the Beatific Vision. It is imperative to remember that this is not some mystical experience, but rather a form of gnosis, that is the unmediated knowledge of God.

Hence, esoterism had to move on from paganism to Christianity. Although necessary from the perspective of initiation, the transition was not without its difficulties exoterically\footnote{\url{https://gornahoor.net/?p=11451}}. The transition disrupted social institutions and was marred by frequent and bloody wars.

\paragraph{Dante}


Dante acknowledges his connection to Virgil, who effectively replaces the Cumaean Sibyl while leading him through hell, and even into Purgatory. Virgil is gone before Dante could reach the Earthly Paradise, or the Primordial State. This is a state of supernatural happiness. In that state, God is known indirectly through creation or philosophical reasoning.

Nevertheless, it is still not yet the Beatific Vision. For that, Beatrice has to take over to begin that ascent. It is important not to confound two entirely different notions:

\begin{itemize}


\item \textbf{Metaphysical Intuition}. This is the immediate realization of the truth of metaphysical teachings. It is beyond any particular state of being, and beyond any particular experience. It is beyond nature.

\item \textbf{Mystical Experience}. This is an experience of a state of being higher than the human state. Although such experiences may be rapturous, ecstatic, extraordinary, insightful, and peaceful, they still belong to the phenomenal realm. This is what distinguishes such an experience from metaphysical intuition.
\end{itemize}


For example, Joan of Arc had mystical experiences. Dante, on the other hand, is not describing a mystical experience of a vision of planets, saints, angels, and so on. Rather, he is describing inner states of being and how to progress from one to the next.

\paragraph{Templars}


According to Rene Guenon, Hermetism was certainly known among the Templars, or the Order of the Temple, and included doctrines of Arabic origin. Dante himself knew of these. In particular, he knew of similar journeys through hell into Heaven, e.g., the nocturnal journey of Muhammed. Guenon makes this observation:

\begin{quotex}
There was a time when there were not only hostile relations, as believed by those who stick to appearances, but also active intellectual exchanges between the East and the West, exchanges which took place mainly by means of Orders e.g., the Templars.
\end{quotex}


As we saw with the transition from pagan esoterism to Christian esoterism, this intellectual exchange came at a time of outer turmoil. Valentin Tomberg says that Saint Bernard's exhortation to the Templars needs to be understood like Krishna's exhortation to Arjuna. Arjuna reluctantly had to battle his cousins, but there was a higher purpose behind the appearances. In his words,

\begin{quotex}
Nowadays many criticise the saint for his intervention sanctioning and encouraging the Crusade, but what he did was simply to make an appeal to ``Christian Arjunas'' on the new field of Kurukshetra, where the two armies of Islam and Christianity had already been assembled for a battle without mercy some centuries before him.
\end{quotex}


\paragraph{The Destruction of the Templars}


The politically motivated destruction or the Templars led to the diminution of esoteric knowledge in the West. Those who maintained such knowledge were forced to do so clandestinely. Thus the alchemists had to disguise their true purpose. The Rosicrucians, who also continued the tradition, had to hide behind the story that they had all scattered to the East. That describes their state of consciousness, not their physical locale.

Guenon claimed to trace some evidence of esoteric teaching in Protestantism, although even Jacob Boehme ran afoul of the authorities. From there, the trail went dark, apart from occasional clues. Nevertheless, the Tradition can never disappear; there will always be a remnant.

\paragraph{Recent Developments}


Since Guenon's death, there have been other developments, of which he was unaware, that alter the picture. First of all, there is the remarkable case of Vladimir Solovyov\footnote{\url{https://gornahoor.net/?p=9927}} who had his own realizations of Divine Sophia.

Vladimir Solovyov, in his trips to Egypt, came to the realization that Hermetism supported the theology of the early church fathers. Besides Hermetism, Solovyov claimed to be following in the tradition of Neoplatonism, Cabala, Swedenborg, and Boehme. However, he never did acknowledge any particular training.

Guenon claimed that the full recovery of esoterism in the West will depend on the reunification of the esoteric Templar teachings with exoteric teachings. Movements have been made in that direction. Boris Mouravieff has revealed the Hermetic roots of the esoteric teachings of Eastern Orthodoxy. For the Roman Church, that task was started more explicitly by Valentin Tomberg. He showed to way to reconcile the Church of Peter with the Church of John\footnote{\url{https://gornahoor.net/?p=16881}}.

However, esoteric teaching can no longer be restricted to secret groups guarded by initiations. Initiatic groups replace ``organic growth'' with ``making''; they depend on human effort rather than grace. Tomberg expands on this notion:

\begin{quotex}
the only truly morally founded reason for keeping esotericism ``esoteric'', i.e. for not bringing it to the broad light of day and popularising it, is the danger of the great misunderstanding of confusing the tower of Babel with the tree of Life , as a consequence of which ``masons'' will be recruited instead of ``gardeners''.

The Church was always conscious of this danger. This is why it always insisted --- whilst appreciating and encouraging effort as such --- on the principle of grace as the sole source for advancing on the way of perfection. This is also why it was always suspicious of so-called ``initiation fraternities'' or such groups who formed themselves at its periphery or beyond it. For, leaving rivalry and other human imperfections out of account, the serious reason for the Church to take a negative attitude towards initiation fraternities is the danger of the substitution of building for growth, of ``doing'' for grace, and of ways of specialisation for the way of salvation.
\end{quotex}

\paragraph{Appendix on Rites}


Guenon has claimed that the Latin rites are no longer effective. That is an absurdity, as long as the apostolic succession remains intact. The efficacy of the rites does not depend at all on the character of knowledge of the priest, just his intention.

In a letter to L. C. d'Amiens from May 1935, he hedges a bit asserting that ``the effects of Catholic rites are far from negligible''. Guenon responds to L C, apparently answering a question:

\begin{quotex}
The Masons help themselves to the influence of the Catholic rites, as you say, I certainly do not see the slightest problem.
\end{quotex}


That would make no sense if the Catholic rites are not effective. Hence, they can be followed if no alternative is available. Obviously, being a Muslim in Europe today is much easier than it was in 1935. In my hometown, for example, there are two mosques within a short bike ride. But, unlike Buddhists or Yogis, they don't offer Sufi initiations to the general public.

Guenon never tires of emphasizing the necessity for observing the rites:

\begin{quotex}
For a ``realization'' of an esoteric order, we must not forget that the observance of the rites constitutes the necessary basis here; and it is also obvious that whoever wants the ``more'' must first, and as a precondition, do the ``less'' (i.e., observe the rites that are common to all).
\end{quotex}


In other words, realization requires following an exoteric tradition, even more assiduously than the commoners. Hence, there can be no ``book learning only'' esoterists. One must pray, meditate, participate in the rites or sacraments. Valentin Tomberg makes the same claim.


\flrightit{Posted on 2020-12-28 by Cologero}

\begin{center}* * *\end{center}

\begin{footnotesize}
\begin{sffamily}
\texttt{Michael M on 2020-12-28 at 12:49 said:}

Thank you for the article, summarizes many important points of this project in a clear way.

I wanted to ask a question based on these two statements: ``It is beyond any particular state of being'' \& ``he is describing inner states of being and how to progress from one to the next.''

Is this to say that Dante is describing the inner states of being themselves and how to progress, but not the specific metaphysical intuition, i.e. because this is beyond any communication or realm of experience, to put it in any terms of experience would make it dualistic and not beyond that? To borrow: ``not made out of any kind of substance which needs maintenance from the point of view of the perceiver, but they just happen, as expression of silence and expression of emptiness''.

Non-dualistic intuition then would be in contrast with mystical experiences, as they would be normally described in terms appropriate to the experiencer and what is experienced, i.e. Saints and Angels as St. Joan of Arc testified, not to speak of whatever inner state of being she was `in' so to speak.

\hfill

\texttt{dionysos areopagite on 2020-12-28 at 19:50 said:}

Thank you so much for this incredible post. I wanted to ask about what you think about the thesis that Catholicism is wholly exoteric and esoteric.. that it does not distinguish between them because it completely combines them.. I think this is a thesis of some Catholics who follow Guenon..

\hfill

\texttt{Cologero on 2020-12-29 at 10:23 said:}

First, I don't understand what the thesis even means. Second, not to diminish the man's accomplishments, Guenon is not Procrustes, so we are not compelled to fit everything to suit him. What Guenon writes in terms of principles is one thing, but his opinions are not dogmatic. I think this post made the distinction clear. If you read Dante as poetry or theology, that is exoteric. If you understand him as describing inner spiritual states, that is esoteric. It is unlikely that anyone can make that transition without help and support of some kind from someone who has.

\hfill

\texttt{dionysos areopagite on 2020-12-29 at 15:24 said:}

``First, I don't understand what the thesis even means. Second, not to diminish the man's accomplishments, Guenon is not Procrustes, so we are not compelled to fit everything to suit him. What Guenon writes in terms of principles is one thing, but his opinions are not dogmatic. I think this post made the distinction clear. If you read Dante as poetry or theology, that is exoteric. If you understand him as describing inner spiritual states, that is esoteric. It is unlikely that anyone can make that transition without help and support of some kind from someone who has.''

I think I would completely agree on both points.. I was thinking specifically about Borella who seems to me to bring up valid points about Guenon in his book Christ: The Original Mystery.. He specifically mentions (Matthew 27:51; Mark 15:38; Luke 23:45), the veil being torn asunder, as evidence of Christ as the fully exoteric Revelation of the esoteric, in other words, full harmony between esoteric and exoteric. I thought this was a compelling thesis.. But I am not inclined to wholly agree with him in all his points of criticism of Guenon..

I would mention organizations like Opus Dei as a kind of example of combining the esoteric and exoteric aspects in Catholicism.. Perhaps at first it wouldn't appear so... I don't know how much you know about that organization.. But they are very interesting to me..

I also wanted to ask a further question which is what do you regard as the purpose of the distinction between esoteric and exoteric teaching.. for example, someone like Leo Strauss would probably say that the purpose of esoteric versus exoteric teaching is to maintain the harmony between the elites and the commoners.. I don't know how much the Guenonian perspective differs from this..

Perhaps you have written on this before so I of course do not intend to burden you if you'd like to point me to a previous post... But I am interested because I love this blog and respect your work very much..

\hfill

\texttt{Cologero on 2020-12-29 at 23:07 said:}

Let me start with the basics. Due to health issues, I am extremely vulnerable to Covid; hence, I am living like an anchorite. My Priest calls me monthly to see if I'm ready for a visit. He asks me: are you praying the rosary?, are you making spiritual communions?, etc. NOT are you reading Borella?, what do you think of Dante?

We are called to know, love and serve God; these usually require different interests and abilities. For most people, going to Mass and receiving the sacraments are all they need. Some of us are more inclined to the ``know'' part, hence the interest in the so-called esoteric. In Man and his Becoming, Guenon writes:

\begin{quotex}
Exoterism and esoterism, regarded not as two distinct and more or less opposed doctrines, which would be quite an erroneous view
\end{quotex}

So esoterism is NOT a secret teaching, but a deeper understanding of the exoteric. Why the interest? Because some of us are called in that direction. Why again? because it helps us understand better the teachings of the Fathers, Saints, Mystics, etc. Pace Borella, Jesus taught in parables and there were many things left out of the Gospels.

Guenon may be somewhat inconsistent in his definitions, so I'll stand by the one I quoted. Rather than argue with any of those fellows, I'll give a concrete example. You decide if it makes sense to you.

Metaphysical realization literally means making the metaphysical teaching, which are initially virtual, ``real''. That is a more fruitful starting point than debating esoteric-exoteric. Guenon emphasizes that such a realization requires concentration and meditation --- skills that must be learned from someone else, apart from rare exceptions. Reading another book won't help, nor will intellectual debates, which is a lower form of knowledge than intuition. That is basic theology. We want to reach intuition.

EXAMPLE: The soul consists of a vegetative and an animal soul; those are both carnal and natural. Humans also have an intellectual soul, which is supernatural. You can read that and ``know'' it as a philosophical or metaphysical viewpoint.

However, there is a deeper way. You can know it through direct perception. Through training, you can learn to ``see'' how the three parts of the soul relate and interact with each other. As Paul reminds us, we are at war with ourselves. So this knowledge will help us develop a unified self, the more we learn about our inner states. ``Know thyself'' needs to be taken literally, and not in a trivial sense like ``I don't like broccoli.''

Maybe, you will begin to notice higher states. I have written posts on Dante's precursors like Augustine, Saint Bernard, the Victorines, who have described such states.

I think you can see how debate and quoting ``authorities'' gets you nowhere, just endless discussion and banter. There is a better way.

\hfill

\texttt{dionysos areopagite on 2020-12-31 at 10:57 said:}

Very much appreciate your responses and all your writings, sir. Being given the gift of metaphysical realization or mystical experience is something very precious and rare indeed, and something to be very grateful for.. and has made me personally very interested in the question of the spiritual renewal of our religion and culture.. many people seem to live as if this isn't possible.. Or they feel some kind of desperation in regard the cultural situation.. It's very hard indeed to not feel this way.. Its hard to not to look at people in Church and find them lacking.. It is very hard to consider how God's plan can give them waters to soothe their thirst.. It certainly seems like a more than human task.. but thinking about it alone helps no one.. and so I also have zero interest in endless debate and citing authorities... I seek out conversation with people who are familiar with ideas and authors in order to seriously think through how genuine spiritual renewal is possible.. I think it is only possible as a cooperation between human and divine authority... So I see writings like your blog as possible signs of genuine interest in spiritual renewal.. So again, I thank you.. God bless and Merry Christmas..

\hfill

\texttt{Max on 2021-01-17 at 14:45 said:}

On the question of esoterism in Guénon and Strauss. The latter is concerned with exoteric writing and politics. The aim of the social order is to allow for its members to develop their full human potential. Guénon makes the distinction between lesser and greater mysteries, and says that the ``lesser mysteries effectively stop at the limits of human possibilities, and constitute in respect of the latter only a preparatory stage and are not themselves their own proper end''. In other words, politics does not contain its own end, and therefore no secular politics is ultimately tenable. This is incidentally also the main point at which Strauss criticizes modernity, and is thus in agreement with Guénon.

Strauss is moreover critical of ``historicism'', while Guénon says that the ultimate goal of the lesser mysteries is beyond the temporal condition, where ``freed from time, the apparent succession of things is transmuted into simultaneity''. That means reaching the primordial state, or earthly paradise. As a Platonist, Strauss rejects the notion that time is a condition for knowledge, or the progressive one according to which ``progressives'' are better at understanding the world than anybody else. Although functioning as a guide for some of the more shadowy lower reaches, Strauss does not seem to be aboard the ``move from paganism to Christianity'', as this piece articulates above. He was very much concerned with the social disruptions that followed though, and which still appears unresolved. In principle, he does not like to discuss the Bible, but the reason is not because it has been inconsequential for political philosophy.

I am not aware that Strauss treated of esoteric teaching in writing. He said that there is an outside, not necessarily what the inner meaning of the fact of there being an outside is. It quickly turns into a scandal for literalists that not everything is of surface level, believing it some kind of malevolent conspiracy. There are two conflicting reactions upon confrontation with the state of affiars of an interior/exterior divide. Some would seek to abolish it, and in effect make everything into mere surface, while realists instead accepts it as the nature of reality as we know it. The first practices secrecy out of lower passions such as fear, while the latter conceives of a kind of hiddenness in the world itself, that no human convention can ever overcome, but that nevertheless serves a secret purpose all its own.

Strauss naturally focuses his attention on the phenomenon of exoterism, while his readers are mostly obsessed with esoterism. He did not have a high opinion of people who believe that they understand an author better than he understood himself though. The difficult point to get across is that you will not find the esoteric teaching in a book, but perhaps pointers in the right direction. These are the options, as valid for Nietzsche as for us:

``...Nietzsche could choose one of two ways: he could insist on the strictly esoteric character of the theoretical analysis of life -- that is, restore the Platonic notion of the noble delusion -- or else he could deny the possibility of theory proper and so conceive of thought as essentially subservient to, or dependent on, life or fate...'' -Leo Strauss

So in conclusion, either interiority is real in its own right, or it is merely an epiphenomenon of ``life'' or what have you. A public teaching accomodates the popular opinion, or the particular views of those it aims to reach. It might make use of the prevailing myths of the times, for example those of liberalism in our day. In writing then, ideas are perhaps expressed dialectically or polemically without demonstrating the reasons why, or it provides principles and definitions without drawing out the logical conclusions. The general idea is that something is always and necessarily left out. Good writers does that deliberately, while for the bad, it just happens accidentally. Silence does not speaks when it is merely an accident rather than concentrated.

\hfill

\end{sffamily}
\end{footnotesize}

